%% kapa: exam=1 round=34 year=2023 period=2 form=A subject=appraisal_law subject_ko=감정평가관계법규 seq=8 number=8
%% source: Q-Net 감정평가사 기출문제 (cst00309) artlSeq=5211789 file=raw/1008/5211789/2023년도 제34회 감정평가사 1차시험 2교시.pdf
%% conversion: kordoc 4.24.51 -> tools/gen1.py (auto); review=proofread-sonnet; flags=none
\begin{question}{8}
국토의 계획 및 이용에 관한 법령상 지구단위계획에 관한 설명으로 옳지 않은 것은?
\begin{choices}
  \item 도시지역 내 지구단위계획구역의 지정목적이 한옥마을을 보존하고자 하는 경우에는 지구단위계획으로 「주차장법」 제19조제3항의 규정에 의한 주차장 설치기준을 100퍼센트까지 완화하여 적용할 수 있다.
  \item 주민이 입안을 제안한 경우, 지구단위계획에 관한 도시·군관리계획결정의 고시일부터 3년 이내에 허가를 받아 사업이나 공사에 착수하지 아니하면 그 3년이 된 날에 지구단위계획구역의 지정에 관한 도시·군관리계획결정은 효력을 잃는다.
  \item 도시지역 외에서 지정되는 지구단위계획구역에서는 지구단위계획으로 당해 용도지역 또는 개발진흥지구에 적용되는 건폐율의 150퍼센트 및 용적률의 200퍼센트 이내에서 건폐율 및 용적률을 완화하여 적용할 수 있다.
  \item 도시지역에 개발진흥지구를 지정하고 당해 지구를 지구단위계획구역으로 지정한 경우에는 지구단위계획으로 당해 용도지역에 적용되는 용적률의 120퍼센트 이내에서 용적률을 완화하여 적용할 수 있다.
  \item 국토교통부장관은 「도시개발법」에 따라 지정된 도시개발구역의 전부 또는 일부에 대하여 지구단위계획구역을 지정할 수 있다.
\end{choices}
\end{question}
\answer{2}
